\section{Elementary tools}\label{sec:tools}

\begin{lemma}\label{lower:lem}
Every \term{goldbach-number} is at least $4$.
\end{lemma}

\begin{proof}
If $n=p+q$ with $p,q\in\Prset$ then $p\geq2$ and $q\geq2$, so $n\geq4$.
\end{proof}

The next lemma is the only tool the computations use.  It replaces the search
over ordered pairs by a single set intersection, so that a sieve computed once
answers the question for every $n$ in range.

\begin{lemma}\label{reflect:lem}
For $n\geq4$, the integer $n$ is a \term{goldbach-number} if and only if
\[
  \Prset\cap(n-\Prset)\ \neq\ \emptyset ,
\]
where $n-\Prset=\{n-p:p\in\Prset,\ p\leq n\}$.  Moreover
$\gc(n)=|\Prset\cap(n-\Prset)|$.
\end{lemma}

\begin{proof}
If $(p,q)$ is a \term{goldbach-pair} for $n$ then $p\in\Prset$ and
$p=n-q\in n-\Prset$, so $p$ lies in the intersection.  Conversely, if
$p\in\Prset\cap(n-\Prset)$ then $p=n-q$ for some prime $q\leq n$, and $(p,q)$
is a Goldbach pair.  The two constructions are mutually inverse and preserve
the first coordinate, so they biject the Goldbach pairs for $n$ with the
elements of the intersection, giving the count.
\end{proof}

\begin{cor}\label{even:cor}
Conjecture~\ref{goldbach:conj} holds for $n$ if and only if some prime
$p\leq n-2$ has $n-p$ prime.
\end{cor}

\begin{proof}
Immediate from Lemma~\ref{reflect:lem} together with $q\geq2$.
\end{proof}
